\section{Conclusion}

In conclusion, this paper has systematically addressed the complexities associated with thermally-induced spontaneous demixing and embrittlement in materials. Parameters such as stable spinodal compositions, elastic moduli, diffusion lengths, and embrittlement criteria have been critically examined to deepen our comprehension of material behavior and stability. The methodology presented herein is not only robust but also scalable, making it an invaluable asset for high-throughput datasets, which are increasingly prevalent in materials science and engineering research.

The experimental findings discussed in this work have shed light on the intricacies of plastic embrittlement and its interplay with quantum elastodynamics. The observed phenomena of thermally-induced spontaneous demixing, coupled with the analyses of elastic anisotropy and material stiffening, have significantly contributed to our understanding of the mechanical performance of materials. These insights are instrumental in guiding future material design and optimizing industrial processes, with the potential to revolutionize the field of materials science. The implications of this research are far-reaching, promising enhanced predictive capabilities and tailored material properties for advanced engineering applications.